Pour analyser la syntaxe de la suite de mots obtenue par l'analyse lexicale, on utilise un analyseur syntaxique. L'analyse syntaxique consiste à vérifier si la suite de mots respecte les règles de grammaire du langage. Cette grammaire spécifie donc la syntaxe (priorités de calculs, associativité, etc.) du langage. L'analyse syntaxique est souvent réalisée à l'aide d'arbres syntaxiques abstraits (AST) qui représentent la structure hiérarchique de la suite de mots.\\\par

On utilise des analyseurs syntaxiques comme \code{bison} ou \code{yacc} pour générer automatiquement le code de l'analyseur syntaxique à partir de la grammaire du langage. 
